% Secao 4: A Jornada dos 3 Anos (Matriz Curricular ano a ano)

% Slide 8: Visao Geral da Jornada Formativa
\begin{frame}{A Jornada dos 3 Anos: Do Zero ao Especialista}{\faRoute\ Um percurso passo a passo planejado para quem nunca mexeu em servidores}
  \begin{columns}[t]
    \begin{column}{0.32\textwidth}
      \begin{greencard}{\faSeedling\ 1º Ano: Fundações}
        \tiny
        \textbf{Descobrindo o Mundo Tech}
        \begin{itemize}\setlength{\itemsep}{0.12em}
          \item Como computadores conversam.
          \item Dominando Linux e Windows.
          \item Eletricidade básica e hardware.
          \item Base forte de matérias do Médio.
        \end{itemize}
        \vspace{0.06cm}
        \centering\greenbadge{Nível: Iniciante Curioso}
      \end{greencard}
    \end{column}
    
    \begin{column}{0.32\textwidth}
      \begin{bluecard}{\faTools\ 2º Ano: Mão na Massa}
        \tiny
        \textbf{Construindo a Infraestrutura}
        \begin{itemize}\setlength{\itemsep}{0.12em}
          \item Cabeamento estruturado e racks.
          \item Configuração prática de Servidores.
          \item Lógica de Programação em Python.
          \item Desmontagem e conserto de PCs.
        \end{itemize}
        \vspace{0.06cm}
        \centering\bluebadge{Nível: Técnico Prático}
      \end{bluecard}
    \end{column}
    
    \begin{column}{0.32\textwidth}
      \begin{purplecard}{\faLock\ 3º Ano: Elite \& Nuvem}
        \tiny
        \textbf{Segurança, Wi-Fi e Web}
        \begin{itemize}\setlength{\itemsep}{0.12em}
          \item Cibersegurança e Firewalls.
          \item Redes Sem Fio e antenas de rádio.
          \item Páginas Web e Banco de Dados.
          \item Prática Profissional e portfólio.
        \end{itemize}
        \vspace{0.06cm}
        \centering\purplebadge{Nível: Especialista Pronto}
      \end{purplecard}
    \end{column}
  \end{columns}
  
  \vspace{0.08cm}
  \begin{statcard}
    \tiny \textcolor{IFCEDarkGreen}{\faLightbulb\ \textbf{Pergunta Frequente:} ``Preciso saber programar ou consertar computador para entrar?'' $\rightarrow$ \textbf{Não!} Você aprende tudo do zero com os professores!}
  \end{statcard}
\end{frame}

% Slide 9: 1º Ano
\begin{frame}{1º Ano: O Ponto de Partida no Universo Tech}{\faCompass\ Os primeiros passos na transição do Ensino Fundamental para o Médio Técnico}
  \begin{columns}[t]
    \begin{column}{0.50\textwidth}
      \begin{greencard}{\faMicrochip\ Disciplinas Técnicas do 1º Ano (280h)}
        \tiny
        \begin{itemize}\setlength{\itemsep}{0.15em}
          \item \textbf{Comunicação de Dados e Redes (120h):} Como dados viajam por fios e ondas; modelos de rede, pacotes e endereçamento IP.
          \item \textbf{Sistemas Operacionais (80h):} O que está por trás do Windows e a mágica do terminal Linux; arquivos e processos.
          \item \textbf{Eletricidade Básica (80h):} Circuitos elétricos, medições com multímetro, fontes de alimentação e segurança.
          \item \textbf{Introdução ao Curso e Orientação (40h):} Carreira em TI, ética e rotina de estudos no ensino federal.
        \end{itemize}
      \end{greencard}
    \end{column}
    
    \begin{column}{0.48\textwidth}
      \begin{moderncard}{\faGraduationCap\ Ensino Médio de Excelência (600h)}
        \tiny
        \begin{itemize}\setlength{\itemsep}{0.15em}
          \item \textbf{Língua Portuguesa \& Redação (80h):} Domínio da escrita e interpretação crítica de textos.
          \item \textbf{Matemática (120h):} Álgebra, geometria e raciocínio lógico.
          \item \textbf{Ciências da Natureza (120h):} Física (40h), Química (40h) e Biologia (40h) com práticas experimentais.
          \item \textbf{Ciências Humanas (160h):} História, Geografia, Filosofia e Sociologia.
          \item \textbf{Linguagens:} Espanhol (40h), Artes (40h) e Ed. Física (40h).
        \end{itemize}
      \end{moderncard}
    \end{column}
  \end{columns}
  
  \vspace{0.06cm}
  \begin{center}
    \tiny \textcolor{TechGray}{\faInfoCircle\ Carga Horária Total do 1º Ano: 960h letivas para uma adaptação acolhedora e estimulante.}
  \end{center}
\end{frame}

% Slide 10: 2º Ano
\begin{frame}{2º Ano: Subindo de Nível (Infra, Servidores \& Código)}{\faCubes\ O ano da transformação: muito laboratório, hardware e as primeiras linhas de código}
  \begin{columns}[t]
    \begin{column}{0.50\textwidth}
      \begin{bluecard}{\faServer\ Disciplinas Técnicas do 2º Ano (360h)}
        \tiny
        \begin{itemize}\setlength{\itemsep}{0.15em}
          \item \textbf{Infraestrutura e Projetos de Redes (80h):} Cabeamento estruturado, conexões de fibra, patch panels e racks.
          \item \textbf{Servidores (120h):} Instalação e administração de servidores DNS, DHCP, Web e arquivos em ambientes de rede.
          \item \textbf{Montagem e Manutenção (80h):} Diagnóstico completo de hardware, montagem de PCs e manutenção preventiva/corretiva.
          \item \textbf{Lógica de Programação (80h):} Algoritmos, variáveis, repetições e criação de programas em Python para automação.
        \end{itemize}
      \end{bluecard}
    \end{column}
    
    \begin{column}{0.48\textwidth}
      \begin{greencard}{\faLightbulb\ Inovação \& Ensino Médio (680h)}
        \tiny
        \begin{itemize}\setlength{\itemsep}{0.15em}
          \item \textbf{Projeto Integrador (40h):} Criação de soluções em equipe para problemas da região dos Inhamuns.
          \item \textbf{Empreendedorismo (40h):} Modelos de negócio digitais, startups e precificação de serviços de TI.
          \item \textbf{Língua Inglesa (40h):} O idioma global da documentação tech.
          \item \textbf{Base BNCC (600h):} Português, Matemática e Ciências com foco na preparação para o ENEM.
        \end{itemize}
      \end{greencard}
    \end{column}
  \end{columns}
  
  \vspace{0.06cm}
  \begin{center}
    \tiny \textcolor{NetBlue}{\textbf{\faLaptopCode\ No 2º ano você já é capaz de montar seu próprio computador e configurar uma rede corporativa inteira!}}
  \end{center}
\end{frame}

% Slide 11: 3º Ano
\begin{frame}{3º Ano: A Elite da Conectividade (Segurança, Wi-Fi e Web)}{\faUserShield\ O ciclo de conclusão: tecnologias avançadas, defesa cibernética e preparação para o futuro}
  \begin{columns}[t]
    \begin{column}{0.50\textwidth}
      \begin{purplecard}{\faLock\ Disciplinas Técnicas do 3º Ano (320h)}
        \tiny
        \begin{itemize}\setlength{\itemsep}{0.15em}
          \item \textbf{Gerência e Segurança de Redes (80h):} Firewalls corporativos, criptografia, detecção de invasões e segurança digital.
          \item \textbf{Redes Sem Fio (80h):} Wi-Fi corporativo de alta performance, propagação de sinal, antenas e proteção wireless.
          \item \textbf{Introdução ao Desenvolvimento Web (80h):} Páginas modernas e responsivas em HTML, CSS e JavaScript.
          \item \textbf{Banco de Dados (80h):} Modelagem e consultas SQL para armazenar dados com velocidade e segurança.
        \end{itemize}
      \end{purplecard}
    \end{column}
    
    \begin{column}{0.48\textwidth}
      \begin{redcard}{\faMedal\ Reta Final: Prática \& ENEM (680h)}
        \tiny
        \begin{itemize}\setlength{\itemsep}{0.15em}
          \item \textbf{Prática Profissional Supervisionada (40h):} Vivência real em projetos com supervisão docente.
          \item \textbf{Foco Total no ENEM:} Língua Portuguesa e Redação (120h), Matemática (80h) e Biologia (80h).
          \item \textbf{Disciplinas Optativas:} Oficinas de redação e preparatórios para Olimpíadas de Física e Astronomia.
          \item \textbf{Formatura Dupla:} Conclusão do Médio com registro profissional de Técnico Federal!
        \end{itemize}
      \end{redcard}
    \end{column}
  \end{columns}
  
  \vspace{0.06cm}
  \begin{center}
    \tiny \textcolor{IFCEDarkGreen}{\textbf{\faTrophy\ Resultado Final: 3.000 horas de capacitação integral que transformam você em um profissional tech disputado pelo mercado!}}
  \end{center}
\end{frame}
